\section{Second Pass: Proofs as Programs With Human Readers}
\chapterline

\subsection{The Curry-Howard breakthrough: proofs as executable code}

In 1934, the American logician Haskell Curry noticed an eerie, beautiful parallel: the formal deduction rules of intuitionistic logic mirrored the type signatures of combinators in combinatory logic. Thirty-five years later, in 1969, William Alvin Howard circulated a legendary set of typewritten notes that turned Curry's observation into a deep, exact isomorphism: the \textbf{Curry-Howard correspondence}.

Howard proved that logic and computation are two dialects of the exact same reality:
\begin{itemize}
\item A mathematical proposition is a \textbf{type signature}.
\item A valid mathematical proof is an \textbf{executable program} that implements that type.
\item Proving a theorem is identical to writing a program that successfully type-checks.
\item Simplifying a proof by eliminating detour lemmas (Gentzen's cut elimination) is identical to \textbf{program execution} (evaluating an expression to its return value).
\end{itemize}

An implication $P \implies Q$ is a function taking an argument of type $P$ and returning a value of type $Q$. A universal statement $\forall x \in A, \, P(x)$ is a generic polymorphic function. An existential statement $\exists x \in A, \, P(x)$ is a paired data structure holding a concrete witness alongside verification that the witness satisfies the contract.

\subsection{The human compiler}

Yet there is one crucial distinction between software engineering and pure mathematics: \textbf{a proof does not compile on a silicon chip; it compiles in the mind of a human reader.}

A compiler has infinite patience for cryptic single-letter variable names and arbitrary nested scopes. A human mathematician has limited working memory. To write a proof that convinces, enlightens, and endures, you must manage the reader's cognitive state as carefully as an engineer manages memory registers, variable lifetimes, and modular functional architecture.

\begin{center}
\begin{tikzpicture}[
  >=Latex,
  frame/.style={draw=#1!80!black, fill=#1!10, thick, rounded corners=3pt, font=\small, inner sep=6pt},
  flow/.style={->, thick, draw=gray!70}
]

  % Title
  \node[font=\bfseries, text=structurepurple] at (0, 3.4) {Execution Trace and Lexical Scope of a Proof};

  % Left Column: Active Variable Scope (Memory Stack)
  \node[font=\bfseries, text=sourceblue] at (-3.4, 2.7) {Human Compiler State (Active Scope)};

  \node[frame=sourceblue, text width=5.4cm] (s1) at (-3.4, 1.6) {
    \textbf{Global Environment}:\\
    $\bullet$ Sets $A, B, C$\\
    $\bullet$ Injective maps $f: A \to B$, $g: B \to C$
  };

  \node[frame=targetgreen, text width=5.4cm] (s2) at (-3.4, -0.4) {
    \textbf{Push Local Scope Frame}:\\
    $\bullet$ Fresh arbitrary inputs: $a_1, a_2 \in A$\\
    $\bullet$ Active hypothesis: $g(f(a_1)) = g(f(a_2))$\\
    $\bullet$ Target obligation: $a_1 = a_2$
  };

  \node[frame=quotientorange, text width=5.4cm] (s3) at (-3.4, -2.4) {
    \textbf{Scope Discharged (Return)}:\\
    $\bullet$ Conclusion certified: $a_1 = a_2$\\
    $\bullet$ Pop local scope frame\\
    $\bullet$ Universal guarantee exported to parent
  };

  % Right Column: Proof Text Steps
  \node[font=\bfseries, text=targetgreen] at (3.4, 2.7) {Written Proof Steps (Instructions)};

  \node[frame=sourceblue, text width=5.4cm] (t1) at (3.4, 1.6) {
    \textbf{Step 1 (Setup)}:\\
    ``To prove $g \circ f$ is injective, let $a_1, a_2 \in A$ be arbitrary.''
  };

  \node[frame=targetgreen, text width=5.4cm] (t2) at (3.4, -0.4) {
    \textbf{Steps 2--3 (Transformations)}:\\
    ``Assume $(g \circ f)(a_1) = (g \circ f)(a_2)$.\\
    By injectivity of $g$, $f(a_1) = f(a_2)$.\\
    By injectivity of $f$, $a_1 = a_2$.''
  };

  \node[frame=quotientorange, text width=5.4cm] (t3) at (3.4, -2.4) {
    \textbf{Step 4 (Conclusion)}:\\
    ``Since $a_1 = a_2$, $g \circ f$ is injective.''\\
    \hfill $\square$
  };

  % Horizontal Aligning Connectors
  \draw[flow] (t1) -- (s1);
  \draw[flow] (t2) -- (s2);
  \draw[flow] (t3) -- (s3);

  % Vertical connectors
  \draw[flow] (t1) -- (t2);
  \draw[flow] (t2) -- (t3);
  \draw[flow] (s1) -- (s2);
  \draw[flow] (s2) -- (s3);

\end{tikzpicture}
\end{center}

\subsection{Variable scope: allocation and deallocation}

Every time you write ``Let $x \in S$ be arbitrary'' or ``Choose $k \in \Z$ such that $n = 2k$,'' you are allocating a variable name in the reader's working memory.

\subsubsection{1. Dangling variables}
A variable cannot be referenced outside the lexical block in which it was introduced.
Consider a proof by cases:
\begin{quote}
\emph{Flawed reasoning}:
``Case 1 ($n$ is even): Then $n = 2k$. \\
Case 2 ($n$ is odd): Then $n = 2m + 1$. \\
Therefore $n(n+1) = 2k(2m+1) \dots$''
\end{quote}
The variable $k$ was defined strictly inside Case 1, and $m$ was defined strictly inside Case 2. In any single universe, either $n$ is even or $n$ is odd; they never exist simultaneously! Multiplying $2k$ by $2m+1$ is a catastrophic scope leakage.

\subsubsection{2. Shadowing and premature binding}
If your theorem already contains a fixed parameter $c$, never write:
\begin{quote}
``Let $c$ be an arbitrary element of $S$.''
\end{quote}
Re-using the letter $c$ shadows the original parameter, erasing its identity in the reader's mind. Introduce a fresh symbol instead: ``Let $x \in S$ be arbitrary.''

\subsection{Managing cognitive load with orienting signposts}

A human compiler cannot look ahead three pages to guess why you suddenly introduced an auxiliary polynomial or rearranged an inequality. You must provide orienting signposts before diving into calculation.

Compare these two styles of writing:

\begin{quote}
\textbf{Cryptic style}:
\[
\begin{aligned}
&f(a_1) = f(a_2) \\
&3a_1 + 4 = 3a_2 + 4 \\
&3a_1 = 3a_2 \\
&a_1 = a_2.
\end{aligned}
\]
\end{quote}
This is a naked list of algebraic manipulations. The reader has no idea whether you are proving injectivity, solving for an unknown, or testing a hypothesis.

Now look at the communicative style:

\begin{quote}
\textbf{Communicative style}:
\begin{proof}
We show that $f$ is injective. Let $a_1, a_2 \in \R$ be arbitrary, and assume that $f(a_1) = f(a_2)$. By definition of $f$, this means
\[
3a_1 + 4 = 3a_2 + 4.
\]
Subtracting $4$ from both sides gives $3a_1 = 3a_2$, and dividing by $3$ yields $a_1 = a_2$. Because equal outputs force equal inputs, $f$ is injective.
\end{proof}
\end{quote}

The second version takes only two additional sentences, but it provides complete orientation:
\begin{itemize}
\item It states the \textbf{goal} upfront (``We show that $f$ is injective'').
\item It explicitly \textbf{allocates variables} with their roles (``Let $a_1, a_2 \in \R$ be arbitrary'').
\item It signals the \textbf{antecedent} (``assume that $f(a_1) = f(a_2)$'').
\item It explicitly announces when the \textbf{postcondition} has been satisfied (``Because equal outputs force equal inputs, $f$ is injective'').
\end{itemize}

\subsection{Exhaustive branch coverage}

When a proof splits into cases, your mathematical compiler checks two invariants:
\begin{enumerate}
\item \textbf{Exhaustiveness}: Do the cases cover every possible element of the ambient domain without exception?
\item \textbf{Independence}: Does each branch reach the target conclusion using only the hypotheses active in that specific branch?
\end{enumerate}

If your domain is the integers $\Z$ and you consider positive integers ($n > 0$) and negative integers ($n < 0$), you have an unhandled branch: $n = 0$. In programming, an unhandled branch leads to a runtime crash; in mathematics, it invalidates the proof.

\begin{labbox}
Take any proof you have written previously. Print it out and annotate it with a colored pen:
\begin{itemize}
\item Draw a box around every variable introduction and label its scope lifetime.
\item Highlight every orienting signpost in green.
\item If you find an algebraic line with no explanation of where it came from or where it is going, rewrite it to guide your human reader.
\end{itemize}
\end{labbox}

\begin{historybox}
The realization that proofs are programs grew out of a fierce philosophical war over the foundations of mathematics. In the early twentieth century, the Dutch mathematician L.E.J. Brouwer founded \emph{intuitionism}, rejecting classical non-constructive proofs. Brouwer argued that claiming an object exists without providing an algorithm to construct it is meaningless; he even rejected the universal validity of the Law of the Excluded Middle ($P \lor \neg P$). David Hilbert vehemently resisted, declaring that denying the mathematician the law of excluded middle was like denying the astronomer the telescope. Classical mathematics won the historical battle, but constructive logic found an unexpected second life with the birth of computer science.

In the 1970s, the Swedish philosopher and logician Per Martin-L\"of developed intuitionistic type theory, providing a unified foundation where mathematical propositions are types and constructive proofs are functional programs. This theoretical bridge sparked the modern revolution of interactive theorem provers: Coq, Isabelle/HOL, and Lean. In 2005, Georges Gonthier used Coq to produce the first fully machine-verified proof of the Four Color Theorem, and in 2012 verified the epic 255-page Feit-Thompson Odd Order Theorem in group theory. In 2021, Fields Medalist Peter Scholze challenged the Lean community to verify his and Dustin Clausen's subtle foundational results in condensed mathematics, discovering that interactive proof assistants could help elite researchers verify proofs at the absolute research frontier.

Yet the greatest trap in formalization is forgetting the human audience. Writing a proof purely to satisfy a mechanical type-checker often produces hundreds of lines of cryptic, unreadable code. A computer can parse a fifty-level nested inductive term without breaking a sweat, but human working memory can only track a few active variables at a time. The true art of mathematical writing is creating proofs that satisfy the formal type-checker in the background while telling a luminous, illuminating conceptual story to the human reader.
\end{historybox}

\subsection{Compiling in human minds---and the structural horizon}

Viewing proofs through the lens of program architecture transforms mathematical exposition. Managing lexical scope, declaring variables before referencing them, providing clear orienting signposts, and ensuring exhaustive case coverage eliminates the most common sources of reader confusion. When you treat your reader's mind as a compiler with limited cache, you naturally structure your arguments into modular, digestible lemmas that make truth self-evident.

Yet treating proofs as programs has clear boundaries. Fully formalizing mathematical proofs in interactive assistants like Lean requires enormous effort, often expanding a single paragraph of textbook reasoning into hundreds of lines of explicit formal tactics. More fundamentally, a mechanical proof assistant certifies that a theorem is correct, but it cannot tell you whether the theorem is interesting, what geometric intuition spawned it, or how it connects to other fields. Formality verifies truth; only human insight discovers meaning.

Having established the foundational mechanics of logic, quantifiers, definitions, and proof architecture, we now turn to the concrete structures that mathematics builds. Chapter 31 introduces \textbf{Second Pass: Set Algebra and Boolean Thinking}, bridging set-theoretic operations with Boolean algebra, algebraic symmetries, and Claude Shannon's digital switching revolution.

\newpage
